% TOP_PROBLEM: {"id": 20002501, "title": "A non-polynomial complex potential with entirely real spectrum", "category_id": 16, "category": {"id": 16, "name": "physics", "display_name": "Mathematical Physics", "description": "Problems at the intersection of mathematics and physics.", "slug": "physics", "order_index": 16, "created_at": "2026-07-31T15:26:25.671Z"}, "status": "solved", "statusQualification": "The inherited solved label is not fresh verification that all three source items have been settled. The title corresponds to only the first item, and no proof endorsement is supplied. The polynomial comparison and known basis case are historical source context. The catalogue title/status was paired with a machine-generated research attempt; that attempt is not peer-reviewed evidence of settlement. The statement covers the original source questions rather than only the title’s subcase."}
% FIRST_POSTED: 2026-10-11
% ENTRY_KIND: statement_only
% SUBMISSION: {"kind": "statement_only", "category": "top_problems", "problem_id": "20002501", "proposal_type": "established", "source": "American Institute of Mathematics, Mathematical aspects of physics with non-self-adjoint operators; workshop 8–12 June 2015, Problems 14.1–16.1. https://aimath.org/pastworkshops/nonselfadjointproblems.pdf", "status_evidence": "Inherited source-catalogue status solved preserved. Background has machine-generated research claims explicitly marked not peer reviewed; these are not independent status evidence. Primary source accessed 11 October 2026; compared against the supplied catalogue excerpt. Historical wording and partial-result remarks were checked for scope; no comprehensive current-literature or proof-verification audit was performed.", "status_checked_on": "2026-10-11", "duplicate_check": "Full 15,773-record source registry checked for identical normalized statements; available site registry checked for identical normalized title/statement and canonical-ID collision. Canonical ID is absent from the site registry. No mathematical-equivalence literature audit performed.", "attribution": "Kwang Shin; Petr Siegl; David A. Smith. Statement preparation: OpenAI GPT-6 Codex assistant, generation 11 October 2026; source-matching and notation editing by the assistant, no independent human mathematical verification."}
% SOURCE_STATEMENT: "14. Kwang Shin\n\n14.1. Non-polynomial complex potentials. Consider\n\nH = − d2\n\ndx 2 + xm + a1xm−1 + · · · + am\n\nin L2(R+) with y(0) cos θ + y′(0) sin θ = 0, where aj ∈ C and θ ∈ C. It is known that H has infinitely many eigenvalues. Moreover, all eigenvalues are real if and only if aj ∈ R for all j and θ ∈ R.\n\nOpen Problem: Is there any non-self-adjoint non-polynomial potential case H, either in L2(R) or in L2(R+),\n\nthat generates infinitely many real eigenvalues and at most finitely many non-real eigenvalues? 15. Petr Siegl\n\n15.1. Riesz basis for Schr¨ odinger operator with complex potential. Consider the self-adjoint har-monic oscillator A0:= −d2/dx2 + x2 in L2(R). Define A:= A0 + V with a complex-valued potential\n\nV ∈ L∞(R).\n\nOpen problem: Is the eigensystem of A a Riesz basis? It is known that the answer is 'yes' if V ∈ Lp(R) for some 1 ≤ p < ∞.16. David A. Smith\n\n16.1. Spectral representation of two-point differential operators. Augmented eigenfunctions are a class of spectral functionals which have been shown to be useful in expressing solutions of initial-boundary value problems [FS15, PSss, Smi14]. This is particularly important in the case where the spatial differential operator is degenerate irregular in the sense of Locker [Loc08], as no other effective solution representation is known. They also provide a spectral theorem where the inverse of the operator is diagonalized.\n\nOpen problems: i) Are there other applications for augmented eigenfunctions? ii) Can a spectral theory be developed using augmented eigenfunctions?\n\n> 6\n\n*Other Open Problems"
% Preparation and final self-audit: CONTRIBUTING.md prompts C and D; source/audit artifacts retained in the local physics_statements_2026_10_11 batch.
% Imported from: physics_statements.tex; UnsolvedMath ID 20002501
% !TeX program = lualatex
\documentclass[11pt]{article}
\usepackage[margin=25mm]{geometry}
\usepackage{fontspec}
\setmainfont{Latin Modern Roman}
\usepackage{amsmath,amssymb,mathtools}
\usepackage{unicode-math}
\setmathfont{Latin Modern Math}
\usepackage{xurl,hyperref}
\hypersetup{colorlinks=true,urlcolor=blue}
\setcounter{secnumdepth}{0}
\setlength{\parindent}{0pt}
\setlength{\parskip}{5pt}
\setlength{\emergencystretch}{3em}
\newcommand{\sref}[2]{\hyperlink{source:#1:#2}{[S#2]}}
\newcommand{\cataloguescope}[1]{\paragraph{Recorded target.}#1}
\newcommand{\R}{\mathbb R}
\newcommand{\C}{\mathbb C}
\newcommand{\N}{\mathbb N}
\newcommand{\Z}{\mathbb Z}
\newcommand{\Q}{\mathbb Q}
\newcommand{\F}{\mathbb F}
\newcommand{\eps}{\varepsilon}
\begin{document}
% UnsolvedMath ID 20002501; source catalogue code AIM-PHYSICS-0009
\setcounter{section}{20002500}
\section{A non-polynomial complex potential with entirely real spectrum}\label{20002501}
{\small UnsolvedMath ID 20002501 · Mathematical Physics}
\subsection{Definitions and mathematical statement}

\paragraph{Shared notation.}$\mathbb N=\{1,2,\ldots\}$; the standard number fields have their usual meanings. All additional conventions are stated locally.


\paragraph{Definitions and notation.}A Schrödinger realization is specified by its potential, domain and boundary conditions. On $\mathbb R_+$ a Robin-type boundary condition can be written $y(0)\cos\theta+y'(0)\sin\theta=0$. A Riesz basis is the image of an orthonormal basis under a bounded invertible operator. An eigensystem may involve generalized eigenvectors when the operator is not diagonalizable. Augmented eigenfunctions in the source are spectral functionals used to represent evolution problems.

\paragraph{Catalogue wording retained for traceability.}
The record bundles source items 14.1--16.1.
\begin{enumerate}
\item Find a non-self-adjoint Schrödinger operator on $L^2(\mathbb R)$ or $L^2(\mathbb R_+)$ with a nonpolynomial potential, infinitely many real eigenvalues and at most finitely many non-real eigenvalues. The source compares with $H=-d^2/dx^2+x^m+a_1x^{m-1}+\cdots+a_m$ on $\mathbb R_+$ and the boundary condition $y(0)\cos\theta+y'(0)\sin\theta=0$, $a_j,\theta\in\mathbb C$.
\item Let $A_0=-d^2/dx^2+x^2$ on $L^2(\mathbb R)$ and $A=A_0+V$ for complex $V\in L^\infty(\mathbb R)$. Is its eigensystem a Riesz basis? The source records an affirmative result for $V\in L^p(\mathbb R)$, some $1\le p<\infty$.
\item Find additional applications for augmented eigenfunctions, and develop a spectral theory using them, beyond the two-point differential-operator initial-boundary representations described in the source.
\end{enumerate}
\paragraph{Scope and source qualification.}The inherited solved label is not fresh verification that all three source items have been settled. The title corresponds to only the first item, and no proof endorsement is supplied. The polynomial comparison and known basis case are historical source context. The catalogue title/status was paired with a machine-generated research attempt; that attempt is not peer-reviewed evidence of settlement. The statement covers the original source questions rather than only the title’s subcase.
\subsection{Short English statement}
Find nonpolynomial complex potentials with almost entirely real spectra, determine a harmonic-oscillator basis property, and develop augmented-eigenfunction spectral theory.
\subsection{Sources}
\begin{itemize}
\item[\hypertarget{source:20002501:1}{S1}] ULAM.AI and the UnsolvedMath Contributors, UnsolvedMath, record 20002501 (AIM-PHYSICS-0009). Catalogue curation under CC BY 4.0; definitions and English exposition added locally.
\url{https://huggingface.co/datasets/ulamai/UnsolvedMath}.
\item[\hypertarget{source:20002501:2}{S2}] American Institute of Mathematics, Mathematical aspects of physics with non-self-adjoint operators; workshop 8–12 June 2015, Problems 14.1–16.1.
\url{https://aimath.org/pastworkshops/nonselfadjointproblems.pdf}.
\end{itemize}
\label{end:20002501}
\end{document}
